\chapter{Langage, ensembles et preuves}
\section*{Exercice 1.1 -- Ensembles}
$A=\{2,4,6,8,10\}$, $B=\{3,6,9\}$. L'intersection est $\{6\}$ (cardinal $1$), la réunion $\{2,3,4,6,8,9,10\}$ (cardinal $7$), le complémentaire $\{1,3,5,7,9\}$ (cardinal $5$). Le complémentaire est relatif à $E$.

\section*{Exercice 1.2 -- Couples}
On obtient $(0,a),(0,b),(0,c),(1,a),(1,b),(1,c)$, soit six couples. L'ordre distingue $(0,a)$ et $(a,0)$ ; le second n'appartient même pas à $A\times B$.

\section*{Exercice 1.3 -- Négations}
Les négations sont $x\leq2$ ; « $x<0$ ou $x\geq5$ » ; « $x\ne-1$ et $x\ne1$ ». Pour sortir d'une intersection, il suffit de manquer une condition ; pour sortir d'une réunion, il faut manquer les deux.

\section*{Exercice 1.4 -- Réciproque}
Si $n=6k$ avec $k\in\Z$, alors $n=3(2k)$, donc $3$ divise $n$. La réciproque est fausse : $3$ est multiple de $3$ sans être multiple de $6$.

\section*{Exercice 1.5 -- Une égalité d'ensembles}
Un élément du membre de gauche est dans $A$ et dans au moins l'un de $B,C$ ; il appartient donc à $A\cap B$ ou à $A\cap C$. Réciproquement, appartenir à l'un de ces deux ensembles impose l'appartenance à $A$ et à $B\cup C$. Les deux inclusions sont établies.

\section*{Exercice 1.6 -- L'ordre des mots}
La première est vraie : choisir $y=-x$, qui dépend de $x$. La seconde est fausse : $x=0$ imposerait $y=0$, alors que $x=1$ imposerait $y=-1$. Un même $y$ ne satisfait pas les deux exigences.
